% CP-VI-0090
% target_chapter: VI/36 — Projective Space
% solution_provenance: FRESH_CANONICAL_AUTHORED_SOLUTION

\subsection*{CP-VI-0090: No global affine coordinates on projective space}

Compute $\Gamma(\mathbb P_k^n,\mathcal O)$ for $n\ge1$ and show that every global regular function is constant.

\medskip
\noindent\textbf{Solution.}

Put $S=k[x_0,\ldots,x_n]$ and cover $\mathbb P_k^n$ by
$D_+(x_0),\ldots,D_+(x_n)$. Let $h$ be a global regular function. On
$D_+(x_0)\cong\operatorname{Spec}k[t_1,\ldots,t_n]$, where
$t_i=x_i/x_0$, it is a polynomial $P(t_1,\ldots,t_n)$.

Suppose $P$ is nonconstant and let $d>0$ be its total degree. Its
degree-$d$ homogenization
\[
H(x_0,\ldots,x_n)=x_0^d P(x_1/x_0,\ldots,x_n/x_0)
\]
is homogeneous of degree $d$. Its restriction $H(0,x_1,\ldots,x_n)$
is the nonzero top-degree homogeneous part of $P$. Thus $x_0$ does
not divide $H$. Choose $j\geq1$ for which
$H(0,x_1,\ldots,x_j=1,\ldots,x_n)$ is nonzero. This is possible
because that top-degree polynomial is not zero.

On the chart $D_+(x_j)$ use coordinates $z_i=x_i/x_j$ ($i\ne j$).
The restriction of $h$ from $D_+(x_0)$ to the overlap is
\[
h=\frac{H(z_0,z_1,\ldots,1,\ldots,z_n)}{z_0^d}.
\]
The numerator is not divisible by $z_0$, so this expression is not
in the polynomial coordinate ring of $D_+(x_j)$: it has a genuine
pole along $z_0=0$. This contradicts the regularity of $h$ on
$D_+(x_j)$. Therefore $d=0$ and $P$ is constant.

The chart $D_+(x_0)$ is dense in the integral scheme $\mathbb P_k^n$;
regular functions agreeing there agree globally. Conversely every
constant defines a global regular function. Hence
\[
\boxed{\Gamma(\mathbb P_k^n,\mathcal O_{\mathbb P_k^n})=k}
\]
for every field $k$ and every $n\geq1$.
